\documentclass[12pt]{article}

\usepackage[utf8]{inputenc}
\usepackage[T1]{fontenc}
\usepackage{mathptmx}
\usepackage[margin=1in]{geometry}
\usepackage[hidelinks]{hyperref}

\setlength{\parindent}{0pt}
\setlength{\parskip}{10pt}
\pagestyle{empty}

\begin{document}

Shreyan Kancharla \\
Davidson College, Davidson, NC, USA \\
shkancharla@davidson.edu \\
ORCID: 0009-0009-8654-110X

\bigskip

2026-09-12

\bigskip

Editor-in-Chief \\
\textit{Diagnostic and Prognostic Research}

\bigskip

Dear Editor,

We submit our manuscript, ``Reported biomarker-cognition associations do not translate into predictive utility: a pre-registered analysis of NHANES 2011-2014,'' for consideration as a research article in \textit{Diagnostic and Prognostic Research}.

At least seven published studies report statistically significant associations between blood metals, nutritional biomarkers, and cognitive performance among US adults aged 60 and over in the NHANES 2011-2014 cycles. Every one reports a regression coefficient or an odds ratio; none reports discrimination, calibration, or cross-validated predictive performance, and none benchmarks a biomarker panel against what is already known from demographics alone. We close that gap using the same data and population. Four nested models (demographics; demographics plus blood metals; demographics plus nutritional biomarkers; both blocks) were evaluated by repeated stratified cross-validation, with a pre-registered primary metric --- the change in cross-validated AUC when the biomarker panel is added to demographics --- and a meaningful-difference threshold fixed before any outcome-predictor relationship was examined.

The result is a precise null. The confidence interval for the primary contrast excludes the improvement specified in advance as meaningful, which lets us report that biomarkers are ruled out as conferring a worthwhile predictive gain, rather than merely that we failed to detect one. A sharper, more concrete finding follows from a pre-specified secondary comparison: the full twenty-analyte panel performed significantly worse than a baseline of demographics plus eight items of routine clinical history, consistently across three algorithms. We also show that quantile g-computation --- the mixture-modelling method most commonly used in this literature --- confers no predictive advantage over demographics alone when evaluated on predictive rather than inferential criteria, and that single-run SHAP rankings over these correlated biomarkers are not reproducible across bootstrap resamples.

We believe this manuscript is well suited to \textit{Diagnostic and Prognostic Research} specifically because the journal's remit is the evaluation of prediction models rather than novel exposures or algorithms --- the contribution here is the question, the evaluation framework, and TRIPOD+AI reporting, applied to a literature that has so far reported associations only. We disclose this candidly in the manuscript. A pre-registered, precise null is, we think, exactly the kind of methodologically rigorous negative result the journal exists to publish, and we report it as such rather than reframing it toward a more favourable-sounding claim.

The manuscript is reported in full against the TRIPOD+AI checklist (Additional file 2). All hypotheses, the primary metric, the meaningful-difference threshold, and every sensitivity analysis were fixed in writing, dated, before any outcome-predictor relationship was examined; the pre-registration document is available as Additional file 1 and, on request, the dated record of its timing.

One reference we want to flag rather than have you find unannounced: reference 1 (Wang et al.) is a Research Square preprint, not yet peer reviewed, and is labelled as such both in the reference list and at its first citation in the Background. We understand your submission guidelines restrict citable references to articles, clinical trial registration records and abstracts that have been published or are in press, and we are glad to remove or reposition this citation if your editorial policy does not permit preprints --- we kept it because it is, to our knowledge, the only report of blood-metal associations with cognition in this specific population and cycle range beyond the six peer-reviewed studies also cited, and dropping it silently seemed worse than flagging it for your judgment.

This manuscript is not under consideration elsewhere. The authors declare no competing interests, and this research received no specific funding. We look forward to your consideration and are happy to suggest reviewers if useful.

\bigskip

Sincerely,

\bigskip

\textbf{Shreyan Kancharla}, corresponding author \\
Davidson College, Davidson, NC, USA \\
shkancharla@davidson.edu \\
ORCID: 0009-0009-8654-110X \\
on behalf of the authors: Shreyan Kancharla and Neil Patel, who contributed equally to this work and share first authorship.

\end{document}
